% ==========================================
% BAB IV PART 1 - PENDAHULUAN DAN PERBANDINGAN
% ==========================================
\cleardoublepage
\chapter{ARSITEKTUR SISTEM OTOMASI INSPEKSI KONTAINER}
\label{chap:perancangan}

Bab ini menyajikan artefak rancangan sistem otomasi inspeksi kontainer
terintegrasi. Susunan pembahasan dibuat mengikuti tujuan tugas akhir: kebutuhan
utama rancangan, deskripsi umum sistem, rancangan arsitektur sistem, rancangan
arsitektur data, model alur informasi, dan pembagian fungsi antarkomponen.

\section{Deskripsi Umum Sistem}
\label{sec:desain-sistem-umum}

\subsection{Visi dan Prinsip Desain}
\label{subsec:visi-prinsip}

Sistem otomasi inspeksi kontainer dirancang berdasarkan prinsip
arsitektur yang mengedepankan modularitas, skalabilitas, dan keamanan.
Pendekatan desain mengacu pada pemisahan layanan berdasarkan tanggung jawab
utama sistem agar akuisisi data, pengolahan hasil inspeksi, persistensi data,
dan penyajian informasi dapat dikembangkan secara konsisten serta dihubungkan
melalui antarmuka yang jelas. Dalam kerangka tersebut, modularitas diterapkan
agar setiap komponen memiliki tanggung jawab tunggal dan dapat berevolusi
secara independen tanpa mengganggu komponen lain. Standardisasi antarmuka
digunakan untuk memastikan seluruh interaksi antarkomponen berlangsung melalui
kontrak layanan yang konsisten, sehingga integrasi internal tetap terjaga dan
perluasan ke sistem eksternal dapat dilakukan secara bertahap. Pemisahan jalur
komunikasi diterapkan dengan menggunakan HTTP untuk data terstruktur dan
\textit{WebSocket} untuk penyajian video langsung, sedangkan pemrosesan
tersebar digunakan untuk menempatkan fungsi yang sensitif terhadap latensi di
\textit{edge} dan fungsi yang menuntut konsolidasi data di layanan terpusat.
Selain itu, prinsip keamanan berlapis dan pencatatan riwayat diterapkan untuk
memastikan data inspeksi, status layanan, dan artefak visual tetap terlindungi
dan dapat ditelusuri kembali melalui identitas kontainer.

\subsection{Kebutuhan Utama Rancangan}
\label{subsec:kebutuhan-utama-rancangan}

Kebutuhan utama pada rancangan ini adalah kebutuhan berprioritas tinggi yang
langsung menentukan apakah arsitektur dapat mendukung alur inspeksi kontainer
secara utuh. Kebutuhan ini tidak dibaca sebagai daftar seluruh fitur aplikasi,
tetapi sebagai dasar pemilihan komponen, alur komunikasi, struktur data, dan
mekanisme evaluasi. Rangkuman kebutuhan utama ditunjukkan pada
Tabel~\ref{tbl:kebutuhan_utama_rancangan}.

\begin{table}[ht]
  \centering
  \footnotesize
  \caption[Kebutuhan utama rancangan arsitektur]{Kebutuhan Utama Rancangan Arsitektur}
  \label{tbl:kebutuhan_utama_rancangan}
  \setlength{\tabcolsep}{4pt}
  \renewcommand{\arraystretch}{1.12}
  \begin{tabularx}{\textwidth}{>{\raggedright\arraybackslash}p{0.14\textwidth} >{\raggedright\arraybackslash}X >{\raggedright\arraybackslash}p{0.28\textwidth}}
    \toprule
    \textbf{Kode} & \textbf{Kebutuhan Utama} & \textbf{Implikasi terhadap Rancangan} \\
    \midrule
    KU-01 & Komponen \textit{edge}, layanan API, aplikasi web, dan penyimpanan harus terhubung dalam satu alur inspeksi. & Arsitektur dipisahkan per lapisan, tetapi komunikasi antarlapisan didefinisikan melalui kontrak layanan. \\
    KU-02 & Identitas kontainer harus menjadi jangkar integrasi hasil OCR, pemindaian, dan konteks manifes. & Entitas kontainer ditempatkan sebagai pusat agregasi data inspeksi. \\
    KU-03 & Hasil inspeksi harus membentuk riwayat kontainer yang memuat status, waktu, penanggung jawab, dan bukti visual. & Data terstruktur disimpan pada basis data, sedangkan artefak visual dipisahkan ke penyimpanan objek. \\
    KU-04 & Jalur data transaksional dan video langsung harus dipisahkan sesuai karakter bebannya. & HTTP digunakan untuk data terstruktur, sedangkan \textit{WebSocket} digunakan untuk video beranotasi langsung. \\
    KU-05 & Akses sistem dan aktivitas penting harus dapat dikendalikan dan ditelusuri. & Rancangan memasukkan autentikasi, otorisasi, dan pencatatan aktivitas sebagai kapabilitas inti. \\
    \bottomrule
  \end{tabularx}
\end{table}

\FloatBarrier

\subsection{Kapabilitas Arsitektur yang Diperlukan}
\label{subsec:kapabilitas-arsitektur}

Arsitektur sistem memerlukan sejumlah kapabilitas untuk memenuhi kebutuhan
utama, fungsional, dan nonfungsional yang telah diidentifikasi.
Tabel~\ref{tbl:architecture_capabilities} merangkum kapabilitas yang diperlukan
untuk mendukung operasi sistem.

\begin{table}[ht]
  \centering
  \footnotesize
  \caption{Kapabilitas Arsitektur yang Diperlukan}
  \label{tbl:architecture_capabilities}
  \setlength{\tabcolsep}{4pt}
  \renewcommand{\arraystretch}{1.12}
  \begin{tabularx}{\textwidth}{>{\raggedright\arraybackslash}p{0.25\textwidth} >{\raggedright\arraybackslash}X >{\raggedright\arraybackslash}p{0.25\textwidth}}
    \toprule
    \textbf{Kapabilitas} & \textbf{Indikator Rancangan} & \textbf{Bukti pada Rancangan} \\
    \midrule
    Koordinasi Proses & Alur inspeksi memiliki urutan aktivitas, status, dan pemicu antaraktivitas yang eksplisit. & Diagram aktivitas UML dan alur proses utama. \\
    Penyimpanan Data Terstruktur & Data OCR, hasil pemindaian, manifes, status, dan penanggung jawab tersimpan dalam struktur yang dapat dikaitkan ke identitas kontainer. & Model data konseptual dan entitas kontainer. \\
    Penyimpanan Artefak Visual & Citra bukti dan video tidak ditempatkan sebagai beban utama basis data transaksional. & Pemisahan MongoDB dan penyimpanan objek. \\
    Pertukaran Informasi & Komponen \textit{edge}, API, web, dan penyimpanan berkomunikasi melalui jalur dan kontrak yang dibedakan menurut jenis data. & HTTP untuk data terstruktur dan \textit{WebSocket} untuk video langsung. \\
    Penyajian Informasi & Pengguna dapat membaca status, hasil pemindaian, dan riwayat kontainer dari antarmuka yang mengambil data melalui layanan API. & Halaman detail kontainer dan alur pembacaan hasil. \\
    Analisis Visual & Hasil OCR dan pemindaian visual digunakan sebagai masukan yang dikaitkan kembali ke entitas kontainer. & Alur OCR, \textit{defect scan}, dan propagasi nomor kontainer. \\
    Pengamanan Akses & Akses mesin-ke-mesin dan akses pengguna dipisahkan sesuai konteks penggunaannya. & \textit{API key} untuk \textit{edge} dan autentikasi pengguna pada web. \\
    Pencatatan Aktivitas & Aktivitas penting memiliki konteks waktu, sumber, dan hubungan ke kontainer agar dapat ditelusuri ulang. & Model audit, riwayat kontainer, dan metadata inspeksi. \\
    \bottomrule
  \end{tabularx}
\end{table}

\FloatBarrier

Kapabilitas pada Tabel~\ref{tbl:architecture_capabilities} diturunkan dari analisis kebutuhan untuk mendukung \textbf{kemudahan pemeliharaan} melalui modularitas, \textbf{interoperabilitas} melalui kontrak layanan yang jelas, \textbf{auditabilitas} melalui pencatatan aktivitas, dan \textbf{keandalan operasi}. Pada realisasinya, kapabilitas tersebut diwujudkan melalui pemisahan komponen \textit{edge}, layanan API, aplikasi web, dan lapisan penyimpanan data, sedangkan interoperabilitas dengan sistem eksternal diposisikan sebagai kesiapan perluasan.

\section{Desain Arsitektur Sistem Terintegrasi}
\label{sec:perbandingan-sistem}

Bagian ini menjawab tujuan pertama tugas akhir, yaitu merancang arsitektur
sistem yang mendefinisikan komponen utama, hubungan antarkomponen, dan alur
informasi inspeksi kontainer secara menyeluruh. Untuk memperjelas dasar
perancangan tersebut, bagian ini terlebih dahulu menyajikan perbandingan
konseptual antara sistem inspeksi kontainer
yang berjalan saat ini di pelabuhan dengan sistem terintegrasi yang
diusulkan. Perbandingan ini menunjukkan transformasi fundamental dalam cara
inspeksi kontainer dilakukan, dari pendekatan manual yang terpisah-pisah
menuju ekosistem terintegrasi yang mengikuti standar.

\subsection{Deskripsi Sistem Saat Ini}
\label{subsec:deskripsi-sistem-saat-ini}

Berdasarkan analisis yang dilakukan di Bab III, sistem inspeksi kontainer yang
berjalan saat ini di sebagian besar pelabuhan Indonesia masih menggunakan
pendekatan hibrida yang menggabungkan pemeriksaan manual dengan teknologi sensor
yang beroperasi secara terpisah. Karakteristik utama sistem saat ini meliputi:

\begin{enumerate}
  \item Akuisisi data terpisah, yaitu perangkat inspeksi beroperasi secara
  independen tanpa antarmuka data yang terdokumentasi. Setiap perangkat
  menghasilkan data yang harus dianalisis ulang secara manual oleh operator yang
  berbeda, sehingga waktu konsolidasi bertambah dan risiko inkonsistensi
  meningkat.
  \item Pemeriksaan manual dengan \textit{checklist}, yaitu petugas inspeksi
  melakukan pemeriksaan visual berdasarkan panduan inspeksi standar. Namun,
  penerapan panduan ini sangat bergantung pada interpretasi subjektif setiap
  petugas dan belum didukung mekanisme validasi konsistensi antarpetugas.
  \item Pencatatan data yang tidak mengikuti standar, yaitu hasil inspeksi
  dicatat secara manual dengan format yang berbeda-beda. Variasi nomenklatur dan
  kategori kerusakan antarpetugas mengharuskan normalisasi ulang sebelum analisis
  historis dan pelacakan tren.
  \item Sistem informasi yang tidak terintegrasi, yaitu data inspeksi tersimpan
  di sistem yang terpisah dari sistem operasional eksternal. Kondisi ini
  memerlukan entri ulang data secara manual ketika informasi inspeksi perlu
  dibagikan kepada pihak lain atau sistem eksternal.
  \item Pelaporan dengan jeda administratif, yaitu laporan inspeksi baru tersedia
  setelah konsolidasi data, verifikasi manual, dan proses pelaporan selesai
  dilakukan.
  \item Ketiadaan jejak audit komprehensif, yaitu belum ada pencatatan sistematis
  untuk melacak siapa yang melakukan inspeksi, kapan inspeksi dilakukan, dan
  perubahan apa yang terjadi pada data.
\end{enumerate}

\subsection{Deskripsi Sistem Usulan}
\label{subsec:deskripsi-sistem-usulan}

Sistem otomasi inspeksi kontainer yang diusulkan mentransformasi
seluruh alur inspeksi menjadi proses yang mengikuti standar, terintegrasi, dan
dapat ditelusuri secara menyeluruh. Sistem ini dirancang berdasarkan
arsitektur tersebar yang mengoptimalkan responsivitas, keandalan, dan
skalabilitas.

Karakteristik utama sistem usulan meliputi:

\begin{enumerate}
  \item Akuisisi data terkoordinasi, yaitu perangkat sensor terintegrasi melalui
  komponen koordinasi yang melakukan sinkronisasi pengambilan data.
  \textit{Metadata} lengkap secara otomatis dilampirkan pada setiap data yang
  dikumpulkan, sedangkan komponen lokal melakukan pemrosesan awal untuk
  standardisasi format dan pemeriksaan kualitas.
  \item Deteksi kerusakan berbasis analisis visual, yaitu modul analisis visual
  melakukan deteksi otomatis terhadap berbagai jenis kerusakan berdasarkan model
  analisis yang dilatih. Setiap deteksi disertai tingkat keyakinan dan tingkat
  keparahan agar dapat divalidasi oleh petugas.
  \item Komunikasi layanan standar, yaitu perubahan status inspeksi, hasil
  deteksi, dan aktivitas penting dipertukarkan melalui antarmuka layanan yang
  konsisten untuk menjaga kejelasan kontrak data antarkomponen.
  \item Pelaporan terstruktur, yaitu hasil inspeksi tersedia secara
  otomatis dalam waktu singkat setelah proses inspeksi dimulai, dapat diakses
  oleh pemangku kepentingan yang berwenang, dan dapat disertai notifikasi ketika
  kondisi kritis ditemukan.
  \item Integrasi dengan ekosistem pelabuhan, yaitu sistem dirancang agar siap
  diintegrasikan dengan sistem operasional eksternal melalui antarmuka layanan
  yang mengikuti standar. Pada ruang lingkup implementasi tugas akhir ini,
  realisasi integrasi masih berfokus pada alur internal proyek, sedangkan
  sinkronisasi ke sistem eksternal diposisikan sebagai arah pengembangan
  lanjutan.
  \item Jejak audit dan kepatuhan, yaitu sistem dirancang untuk menjaga riwayat
  transaksi, perubahan data, akses informasi, dan tindakan
  pengguna melalui pencatatan stempel waktu, identitas pengguna, serta artefak
  inspeksi yang terkait dengan kontainer.
  \item Ketahanan operasional lokal, yaitu komponen lokal tetap menjalankan
  akuisisi video, inferensi, dan pemantauan aliran secara dekat dengan sumber
  data, sementara penyimpanan hasil inspeksi tetap mengikuti alur terstruktur
  melalui layanan API saat konektivitas tersedia.
\end{enumerate}

\subsection{Perbandingan Konseptual Sistem}
\label{subsec:perbandingan-komprehensif}

Tabel~\ref{tbl:comparison_characteristics} menyajikan perbandingan karakteristik konseptual antara sistem saat ini dan sistem usulan.

\begin{table}[ht]
  \centering
  \footnotesize
  \renewcommand{\arraystretch}{1.2}
  \setlength{\tabcolsep}{4pt}
  \caption[Perbandingan sistem saat ini dan usulan]{Perbandingan Sistem Saat Ini dan Sistem Usulan}
  \label{tbl:comparison_characteristics}
  \begin{tabularx}{\textwidth}{>{\raggedright\arraybackslash}p{0.16\textwidth} >{\raggedright\arraybackslash}X >{\raggedright\arraybackslash}X >{\raggedright\arraybackslash}p{0.14\textwidth}}
    \toprule
    \textbf{Aspek} & \textbf{Sistem Saat Ini} & \textbf{Sistem Usulan} & \textbf{Dampak} \\
    \midrule
    Akuisisi Data & Perangkat terpisah, tidak terkoordinasi & Terintegrasi melalui komponen koordinasi dengan \textit{metadata} identitas, waktu, kamera, dan artefak & Konsolidasi ulang data manual berkurang \\
    Deteksi Kerusakan & Inspeksi visual manual, subjektif, inkonsisten & Deteksi berbasis analisis otomatis dengan tingkat keyakinan yang dapat divalidasi petugas & Keluaran inspeksi memiliki struktur dan bukti visual yang sama \\
    Pemrosesan Data & Analisis manual oleh operator tunggal & Pemrosesan pada komponen lokal dan terpusat sesuai karakter beban & Jalur kerja lokal dan pusat terpisah menurut fungsi \\
    Pelaporan & Konsolidasi manual, format tidak standar, tertunda & Laporan terstruktur setelah data inspeksi diterima dan disimpan & Jeda konsolidasi laporan manual dapat dikurangi \\
    Integrasi Sistem & Entri ulang manual, tidak ada antarmuka standar & Integrasi berbasis antarmuka layanan dengan kontrak data yang konsisten & Entri ulang dapat dikurangi \\
    Notifikasi & Manual, tidak konsisten, tertunda & Notifikasi berbasis aturan kepada pengguna terotorisasi & Distribusi informasi lebih tertata \\
    Jejak Audit & Catatan manual, mudah hilang, tidak terstruktur & Pencatatan sistematis dengan stempel waktu, pelacakan pengguna, dan artefak bukti & Kesiapan kepatuhan meningkat \\
    Skalabilitas & Terbatas oleh jumlah petugas & Pemisahan komponen agar kapasitas dapat dikembangkan per lapisan & Kapasitas dapat dihitung per jalur inspeksi \\
    Keandalan & Titik kegagalan tunggal, tanpa duplikasi & Pemisahan komponen lokal dan terpusat dengan pemantauan status layanan & Gangguan dapat dilokalisasi per komponen \\
    Keamanan Data & Catatan manual tidak terenkripsi, tanpa kontrol akses & Perlindungan data di setiap lapisan dan kontrol akses berbasis peran & Keamanan lebih terkendali \\
    \bottomrule
  \end{tabularx}
\end{table}

\FloatBarrier
\vspace{-2.6cm}

Perbandingan pada Tabel~\ref{tbl:comparison_characteristics} menunjukkan bahwa transformasi dari sistem manual yang
terpisah-pisah menuju sistem terintegrasi diarahkan untuk memperbaiki
alur pencatatan, konsistensi pembacaan hasil inspeksi, dan kesiapan
riwayat kontainer terhadap standar. Dampak yang ditampilkan pada tabel merupakan
dampak rancangan yang diharapkan dari perubahan arsitektural, bukan hasil
pengukuran operasional penuh. Sistem usulan dirancang untuk mengatasi
keterbatasan mendasar sistem saat ini melalui penerapan prinsip modularitas,
standardisasi, dan integrasi yang konsisten dengan kebutuhan yang telah
diidentifikasi dalam Bab III.

\subsection{Rincian Arsitektur Komponen}
\label{sec:arsitektur-komponen}

Bagian ini menjelaskan desain komponen sistem otomasi inspeksi kontainer, termasuk identifikasi komponen utama, tanggung jawab masing-masing komponen, dan pola interaksi antarkomponen. Desain ini kemudian menjadi dasar realisasi sistem pada proyek pengembangan sistem, sehingga pembahasannya menekankan struktur yang tercermin pada hubungan antara proses inspeksi, data, dan antarmuka pengguna.

Sistem diorganisasikan dalam empat lapisan utama: (1) \textbf{\textit{Presentation Layer}} untuk interaksi pengguna melalui aplikasi web, (2) \textbf{\textit{Application Layer}} untuk logika proses inspeksi, autentikasi, dan orkestrasi data melalui layanan API, (3) \textbf{\textit{Data Layer}} untuk penyimpanan data terstruktur dan artefak visual, dan (4) \textbf{\textit{Edge Layer}} untuk akuisisi aliran video, inferensi, dan pengiriman hasil inspeksi dari lapangan.

\subsubsection{Diagram Konteks Sistem}
\label{subsec:system-context}

\textit{Diagram Konteks Sistem} disusun berdasarkan hasil analisis kebutuhan
yang menunjukkan bahwa sistem inspeksi tidak berdiri sendiri, tetapi berada
dalam jaringan interaksi operasional yang melibatkan petugas lapangan,
pengelola terminal, otoritas, dan sistem informasi lain di lingkungan
pelabuhan. Karena itu, diagram ini tidak hanya dimaksudkan untuk menampilkan
siapa yang berhubungan dengan sistem, tetapi juga untuk menjelaskan mengapa
batas sistem perlu dibedakan antara interaksi yang sudah direalisasikan pada
ruang lingkup tugas akhir ini dan interaksi yang masih berada pada tingkat
arsitektur sasaran.

Penyusunan diagram konteks dilakukan melalui tahapan berikut:

\begin{enumerate}
  \item Mengidentifikasi proses utama inspeksi kontainer dari hasil analisis
  Bab III, yaitu permohonan atau pemicu pemeriksaan, akuisisi data lapangan,
  pengolahan hasil inspeksi, verifikasi temuan, dan penyampaian status inspeksi.
  \item Mengelompokkan pihak yang berinteraksi dengan proses tersebut menjadi
  aktor internal, aktor eksternal, dan sistem eksternal. Petugas inspeksi dan
  operator terminal ditempatkan sebagai aktor operasional, sedangkan Bea dan
  Cukai, perusahaan logistik, pemilik kontainer, dan \textit{shipper}/
  \textit{consignee} ditempatkan sebagai pihak penerima atau pemberi informasi
  inspeksi.
  \item Menentukan jenis pertukaran informasi untuk setiap aktor, seperti data
  permohonan layanan, status inspeksi, bukti visual, notifikasi temuan,
  permintaan verifikasi, dan ringkasan hasil inspeksi.
  \item Menetapkan batas sistem dengan menempatkan Cargovision sebagai pusat
  koordinasi data inspeksi, sementara sistem seperti Inaportnet, TOS/PCS, dan
  API Ditjen Hubla diposisikan sebagai sasaran integrasi eksternal yang
  berkaitan dengan konteks pelabuhan.
\end{enumerate}

\begin{figure}[htbp]
  \centering
  \vspace{0.45em}
  \resizebox{0.97\textwidth}{!}{%
  \begin{tikzpicture}[
      font=\scriptsize,
      system/.style={rectangle, draw=black, rounded corners, align=center, text width=4.95cm, minimum height=1.45cm},
      actor/.style={rectangle, draw=black, align=center, text width=3.18cm, minimum height=0.76cm},
      ext/.style={rectangle, draw=black, dashed, align=center, text width=3.18cm, minimum height=0.76cm},
      line/.style={-{Latex[length=1.8mm]}, thick},
      target/.style={-{Latex[length=1.8mm]}, thick, dashed},
      relmark/.style={inner sep=0pt, font=\tiny\bfseries},
      legendbox/.style={rectangle, draw=black, rounded corners, fill=white, inner sep=4pt},
      dashedline/.style={thick, dashed}
    ]
    \node[system] (cv) at (0,0) {\textbf{Cargovision}\\Sistem otomasi inspeksi\\kontainer\\OCR, pemindaian\\riwayat, bukti visual};

    \node[actor] (inspector) at (-6.0,1.9) {Petugas Inspeksi};
    \node[actor] (operator) at (-6.0,-1.9) {Operator Terminal};
    \node[actor] (web) at (0,-3.0) {Aplikasi Web\\Pengguna Internal};

    \node[ext] (customs) at (6.45,2.55) {Bea dan Cukai};
    \node[ext] (logistics) at (6.45,1.05) {Perusahaan Logistik\\Pemilik Kontainer};
    \node[ext] (shipper) at (6.45,-0.65) {\textit{Shipper}/\textit{Consignee}};
    \node[ext] (tos) at (6.45,-2.35) {TOS / PCS};
    \node[ext] (inaportnet) at (0,2.55) {Inaportnet\\API Ditjen Hubla};

    \draw[line] (inspector.east) -- ([yshift=0.32cm]cv.west);
    \draw[line] (operator.east) -- ([yshift=-0.32cm]cv.west);
    \draw[line] ([xshift=0.42cm]cv.south) -- ([xshift=0.42cm]web.north);
    \draw[line] ([xshift=-0.42cm]web.north) -- ([xshift=-0.42cm]cv.south);
    \node[relmark] at (-3.30,1.86) {1};
    \node[relmark] at (-3.30,-1.86) {2};
    \node[relmark] at (0.78,-1.65) {3};
    \node[relmark] at (-0.78,-1.65) {4};

    \coordinate (externalHub) at (4.55,0);
    \draw[dashedline] (cv.east) -- (externalHub);
    \draw[target] (externalHub) |- (customs.west);
    \draw[target] (externalHub) |- (logistics.west);
    \draw[target] (externalHub) |- (shipper.west);
    \draw[target] (externalHub) |- (tos.west);
    \draw[target] (cv.north) -- (inaportnet.south);
    \node[relmark] at (0.42,1.70) {5};
    \node[relmark] at (4.88,0.42) {6};
    \node[legendbox, anchor=north west, minimum width=4.90cm, minimum height=0.90cm] (legend) at (-6.8,-3.00) {};
    \draw[line] ([xshift=0.22cm,yshift=-0.28cm]legend.north west) -- ++(0.78,0)
      node[anchor=west, xshift=0.14cm] {realisasi internal};
    \draw[target] ([xshift=0.22cm,yshift=-0.62cm]legend.north west) -- ++(0.78,0)
      node[anchor=west, xshift=0.14cm] {sasaran integrasi eksternal};
  \end{tikzpicture}}
  \vspace{0.15em}

  {\scriptsize
  \setlength{\tabcolsep}{3pt}
  \begin{tabular}{@{}cl@{\hspace{1.0em}}cl@{\hspace{1.0em}}cl@{}}
    \textbf{1} & validasi temuan &
    \textbf{2} & status operasional &
    \textbf{3} & hasil/video langsung \\
    \textbf{4} & tinjauan &
    \textbf{5} & integrasi Inaportnet &
    \textbf{6} & integrasi eksternal
  \end{tabular}}
  \caption{Diagram Konteks Sasaran Sistem Otomasi Inspeksi Kontainer}
  \label{fig:system_context}
\end{figure}

\FloatBarrier

Diagram konteks sistem pada Gambar \ref{fig:system_context} menunjukkan
arsitektur sasaran sistem otomasi inspeksi kontainer dalam ekosistem pelabuhan yang
lebih luas. Pada implementasi yang direalisasikan pada tugas akhir ini, batas
integrasi aktual masih berfokus pada interaksi internal antara komponen
\textit{edge}, layanan API, aplikasi web, autentikasi pengguna, dan
penyimpanan artefak visual. Dengan demikian, aktor dan sistem eksternal pada
diagram ini perlu dipahami sebagai konteks sasaran pengembangan, bukan batas
implementasi yang seluruhnya telah direalisasikan.

Aktor manusia yang muncul pada diagram mewakili kebutuhan informasi yang
berbeda. Petugas inspeksi lapangan membutuhkan sarana untuk memverifikasi hasil
deteksi dan menerima notifikasi temuan. Operator terminal memerlukan status
inspeksi untuk mendukung keputusan operasional. Otoritas pemerintah
membutuhkan bukti visual dan hasil deteksi untuk proses kepatuhan, sedangkan
perusahaan logistik, pemilik atau operator kontainer, serta pihak
\textit{shipper}/\textit{consignee} membutuhkan hasil inspeksi untuk
perencanaan tindak lanjut, logistik, dan pengiriman.

Sistem eksternal pada diagram diposisikan sebagai sasaran integrasi, bukan
sebagai komponen yang seluruh hubungannya telah direalisasikan. Sistem
operasional pelabuhan menjadi target untuk menyediakan informasi pergerakan
kontainer dan menerima status inspeksi. Platform layanan pelabuhan diposisikan
sebagai target untuk mengirimkan permohonan layanan yang membutuhkan verifikasi
inspeksi dan menerima kembali status inspeksi. Sementara itu, sistem pemerintah
diposisikan sebagai target untuk sinkronisasi data resmi yang berkaitan dengan
validasi dan pelaporan.

Dengan pembacaan tersebut, batas sistem pada tugas akhir ini tetap jelas.
Realisasi aktual difokuskan pada pengolahan data inspeksi, integrasi internal,
dan pelaporan hasil melalui antarmuka standar, sedangkan
interaksi yang lebih luas pada diagram konteks tetap dibaca sebagai arah
pengembangan arsitektur sasaran.

\subsubsection{Struktur Lapisan Arsitektur}
\label{subsec:arsitektur-komponen-high-level}

Sistem dirancang dengan arsitektur berlapis yang memisahkan tanggung jawab
berdasarkan domain fungsional dan lokasi pemrosesan. Pemisahan ini diturunkan
dari kebutuhan untuk membedakan antarmuka pengguna, orkestrasi layanan,
penyimpanan data, dan pemrosesan lapangan agar masing-masing dapat berkembang
tanpa saling membebani.

Struktur lapisan pada Gambar~\ref{fig:layered_architecture} tidak ditetapkan
secara langsung, tetapi diturunkan dari kapabilitas arsitektur pada
Tabel~\ref{tbl:architecture_capabilities} dan batas interaksi pada
Gambar~\ref{fig:system_context}. Tahap penurunannya adalah sebagai berikut.
Pertama, kebutuhan penyajian informasi kepada operator dan pemangku kepentingan
dikelompokkan sebagai lapisan presentasi. Kedua, kebutuhan koordinasi proses,
autentikasi, pengelolaan inspeksi, dan pertukaran informasi dikelompokkan
sebagai lapisan aplikasi. Ketiga, kebutuhan penyimpanan data terstruktur dan
artefak visual dipisahkan sebagai lapisan data. Keempat, kebutuhan akuisisi
video, inferensi awal, dan komunikasi dari lokasi lapangan ditempatkan sebagai
lapisan \textit{edge}. Hasil pemetaan tersebut diringkas pada
Tabel~\ref{tbl:derivasi-lapisan-arsitektur}.

\begin{table}[ht]
  \centering
  \small
  \caption{Penurunan Lapisan Arsitektur dari Kebutuhan Sistem}
  \label{tbl:derivasi-lapisan-arsitektur}
  \begin{tabular}{p{0.22\textwidth} p{0.31\textwidth} p{0.35\textwidth}}
    \toprule
    \textbf{Lapisan} & \textbf{Dasar Kebutuhan} & \textbf{Tanggung Jawab Arsitektural} \\
    \midrule
    \textit{Presentation Layer} &
    Penyajian informasi dan interaksi pengguna &
    Menampilkan hasil inspeksi, status pemantauan, video langsung, dan akses operasional bagi pengguna. \\
    \textit{Application Layer} &
    Koordinasi proses, autentikasi, dan pertukaran informasi &
    Mengelola logika inspeksi, penerimaan hasil dari \textit{edge}, autentikasi, dan penyediaan antarmuka data. \\
    \textit{Data Layer} &
    Penyimpanan data terstruktur dan artefak visual &
    Menjaga persistensi data inspeksi, data master, dan media bukti agar dapat ditelusuri kembali. \\
    \textit{Edge Layer} &
    Akuisisi data lapangan dan pemrosesan dekat sumber data &
    Menghubungkan kamera, menjalankan inferensi awal, mengelola aliran video, dan mengirim hasil inspeksi. \\
    \bottomrule
  \end{tabular}
\end{table}

\FloatBarrier

Pada \textit{presentation layer}, aplikasi web menyediakan akses bagi operator
untuk memantau aliran video, meninjau hasil identifikasi dan pemindaian,
mengelola aliran kamera, serta membaca hasil inspeksi. Pada
\textit{application layer}, layanan API menerima hasil identifikasi dan
pemindaian dari komponen \textit{edge}, menyimpan data inspeksi, menyediakan
antarmuka data untuk aplikasi web, dan menangani autentikasi. Pada lapisan ini
juga terdapat layanan integrasi aplikasi yang menjaga konsistensi hubungan
antarentitas utama seperti kontainer, manifes, dan hasil pemindaian.

Pada \textit{data layer}, \textit{database} digunakan untuk menyimpan data
transaksional dan data master, sedangkan \textit{file storage} digunakan untuk
menyimpan media inspeksi seperti citra hasil identifikasi, visualisasi deteksi,
dan artefak bukti inspeksi. Adapun \textit{edge layer} beroperasi di lokasi
lapangan untuk akuisisi data dan pemrosesan lokal. Layanan \textit{edge}
menghubungkan kamera \textit{RTSP}, menjalankan inferensi OCR dan deteksi,
mengelola aliran video, dan mengirimkan hasil inspeksi ke layanan API. Di
lapisan yang sama, layanan video langsung menyiarkan \textit{frame}
beranotasi ke aplikasi web melalui \textit{WebSocket} untuk mendukung
pemantauan waktu nyata.

\begin{figure}[htbp]
  \centering
  \resizebox{0.95\textwidth}{!}{%
  \begin{tikzpicture}[
      font=\small,
      layerbox/.style={rectangle, draw=black, rounded corners, minimum width=10.5cm, minimum height=0.95cm, align=center},
      line/.style={-{Latex[length=2mm]}, thick},
      flowlabel/.style={fill=white, inner sep=1.5pt, font=\small}
    ]
    \node[layerbox] (presentation) at (0,0) {\textbf{Lapisan Presentasi}\\Aplikasi Web dan Antarmuka Pengguna};
    \node[layerbox] (application) at (0,-1.45) {\textbf{Lapisan Aplikasi}\\Layanan Inspeksi, Analisis, dan Integrasi};
    \node[layerbox] (data) at (0,-2.9) {\textbf{Lapisan Data}\\Database dan Penyimpanan Berkas};
    \node[layerbox] (local) at (0,-4.35) {\textbf{Lapisan \textit{Edge}}\\Akuisisi Video dan Inferensi Awal};

    \draw[line] ([xshift=-1.25cm]application.north) -- ([xshift=-1.25cm]presentation.south);
    \node[flowlabel] at (-2.1,-0.72) {REST};
    \draw[line] ([xshift=1.35cm]application.south) -- ([xshift=1.35cm]data.north);
    \node[flowlabel] at (2.55,-2.16) {persistensi};
    \draw[line] (local.west) -- ++(-0.85,0) |- node[flowlabel, pos=0.62, left] {hasil inspeksi} (application.west);
    \draw[line, dashed] (local.east) -- ++(1.1,0) |- node[flowlabel, pos=0.74, right] {video langsung} (presentation.east);
  \end{tikzpicture}}
  \caption{Arsitektur Berlapis Konseptual Sistem Otomasi Inspeksi Kontainer}
  \label{fig:layered_architecture}
\end{figure}

\FloatBarrier

Gambar~\ref{fig:layered_architecture} menunjukkan empat lapisan arsitektur
sistem yang diorganisasikan berdasarkan prinsip pemisahan tanggung jawab.
\textbf{\textit{Presentation Layer}} menyediakan antarmuka kepada pengguna melalui
aplikasi web. \textbf{\textit{Application Layer}} mengimplementasikan logika proses inspeksi
utama sistem melalui layanan API yang menerima dan menyajikan data inspeksi.
\textbf{\textit{Data Layer}} menyediakan mekanisme penyimpanan terstruktur dan
penyimpanan media. \textbf{\textit{Edge Layer}} menangani pengumpulan data, inferensi
awal, dan penyiaran video langsung. Diagram ini menggambarkan struktur
konseptual berlapis; pada realisasi aktual, komunikasi utama antarlayanan
dilakukan secara sinkron melalui HTTP untuk data terstruktur dan
\textit{WebSocket} untuk video langsung, tanpa penggunaan \textit{message
queue} sebagai komponen inti.

Pola interaksi utama dijelaskan lebih lanjut pada subbagian berikutnya, meliputi
pola permintaan-respons untuk data terstruktur, pola \textit{WebSocket} untuk
video langsung, dan pola propagasi konteks inspeksi untuk menyelaraskan nomor
kontainer pada beberapa aliran pemindaian.


\subsubsection{Pola Interaksi Komponen}

Interaksi antarkomponen mengikuti pola komunikasi standar untuk menjaga konsistensi dan kemudahan pemeliharaan sebagai berikut:

\begin{enumerate}
  \item Pola permintaan-respons digunakan untuk komunikasi yang memerlukan
  respons langsung dan bersifat transaksional. Pola ini mendukung operasi yang
  memerlukan konfirmasi segera dan menjaga konsistensi transaksi antarkomponen.

  \item Pola video langsung digunakan untuk menyiarkan \textit{frame} beranotasi
  dari komponen \textit{edge} ke aplikasi web secara waktu nyata. Pola ini
  memisahkan kebutuhan pemantauan visual dari jalur data transaksional.

  \item Pola propagasi konteks inspeksi digunakan untuk menyebarkan nomor
  kontainer yang telah dikenali oleh aliran OCR ke aliran pemindaian lain yang
  sedang aktif sehingga hasil beberapa proses inspeksi dapat dikaitkan ke
  entitas kontainer yang sama.
\end{enumerate}

Diagram sekuens UML pada Gambar~\ref{fig:data_flow} diturunkan dari tiga
objek informasi utama yang harus bergerak di dalam sistem, yaitu aliran video
dari kamera, data hasil identifikasi dan pemindaian, serta artefak visual
sebagai bukti inspeksi. Ketiga objek informasi tersebut kemudian dipetakan
terhadap komponen yang bertanggung jawab untuk akuisisi, pemrosesan,
penyimpanan, dan penyajian. Dengan demikian, diagram tidak hanya menggambarkan
koneksi teknis, tetapi juga menunjukkan alasan pemisahan jalur antara data
persisten dan video langsung.

\begin{figure}[htbp]
  \centering
  {\tikzumlset{font=\scriptsize}
  \resizebox{0.98\textwidth}{!}{%
  \begin{tikzpicture}
    \begin{umlseqdiag}
      \umlobject[class={Kamera RTSP}]{Kamera}
      \umlcontrol[class={Cargovision Edge}]{Edge}
      \umlcontrol[class={API Cargovision}]{API}
      \umldatabase[class={Data Kontainer}]{MongoDB}
      \umldatabase[class={Artefak Visual}]{R2}
      \umlboundary[class={Aplikasi Web}]{Web}

      \begin{umlcall}[op={aliranVideo()}, type=asynchron, dt=4]{Kamera}{Edge}
      \end{umlcall}
      \begin{umlcallself}[op={OCR dan validasi nomor}, type=synchron, return={nomor stabil}, dt=5]{Edge}
      \end{umlcallself}
      \begin{umlcall}[op={identifikasiKontainer()}, type=synchron, return={containerId}, dt=4]{Edge}{API}
        \begin{umlcall}[op={upsertContainer()}, type=synchron, return={containerId}]{API}{MongoDB}
        \end{umlcall}
      \end{umlcall}
      \begin{umlcallself}[op={pemindaianKerusakan()}, type=synchron, return={hasil scan}, dt=4]{Edge}
      \end{umlcallself}
      \begin{umlcall}[op={kirimHasilInspeksi()}, type=synchron, return={status tersimpan}, dt=4]{Edge}{API}
        \begin{umlcall}[op={simpanArtefak()}, type=synchron, return={url bukti}]{API}{R2}
        \end{umlcall}
        \begin{umlcall}[op={catatRiwayat()}, type=synchron, return={riwayat diperbarui}, dt=2]{API}{MongoDB}
        \end{umlcall}
      \end{umlcall}
      \umlsdnode[dt=55]{Web}
      \begin{umlcall}[op={ambilDetailKontainer()}, type=synchron, return={riwayat kontainer}, dt=5]{Web}{API}
      \end{umlcall}
      \begin{umlcall}[op={streamVideoAnotasi()}, type=asynchron, dt=4]{Edge}{Web}
      \end{umlcall}
    \end{umlseqdiag}
  \end{tikzpicture}}}
  \caption[Diagram sekuens UML alur inspeksi Cargovision]{Diagram Sekuens UML Alur Inspeksi Cargovision}
  \label{fig:data_flow}
\end{figure}

\FloatBarrier

Gambar \ref{fig:data_flow} mengilustrasikan urutan interaksi utama dalam sistem
otomasi inspeksi kontainer. Alur dimulai dari kamera RTSP yang mengirimkan
aliran video ke komponen \textit{edge}. Komponen \textit{edge} melakukan
inferensi awal untuk membaca nomor kontainer, kemudian layanan API membentuk
atau memperbarui entitas \textit{Container}. Setelah nomor kontainer
terkonfirmasi, hasil pemindaian kerusakan atau kategori menggunakan konteks
kontainer yang sama dan dikirim kembali ke layanan API. Ringkasan hasil
disimpan pada MongoDB, sedangkan artefak visual disimpan pada Cloudflare R2.
Aplikasi web membaca riwayat kontainer melalui layanan API dan menerima video
beranotasi langsung melalui \textit{WebSocket}. Alur ini menunjukkan pemisahan
tanggung jawab antara akuisisi, pemrosesan, penyimpanan, dan penyajian
informasi.

\section{Arsitektur Data Pemeriksaan Kontainer}

Bagian ini menjawab tujuan kedua tugas akhir, yaitu merancang arsitektur data
yang mendukung pencatatan hasil inspeksi, penyimpanan bukti visual, dan
pelacakan riwayat pemeriksaan kontainer. Arsitektur data menjelaskan struktur
informasi yang dikelola sistem dan mekanisme pertukaran data antarkomponen.
Desain data dirancang untuk mendukung kebutuhan fungsional dan nonfungsional
yang telah diidentifikasi dalam Bab III, dengan fokus pada integritas data,
konsistensi, dan kemudahan integrasi.

\subsection{Model Data Konseptual}

Sistem mengelola data menggunakan repositori terstruktur untuk informasi
transaksional dan data master. Model data dirancang dengan prinsip
standardisasi untuk mengurangi duplikasi dan menjaga konsistensi informasi
antarhasil inspeksi. Model ini diturunkan dari kebutuhan integrasi hasil OCR,
hasil pemindaian, dan konteks manifes ke dalam satu entitas kontainer, sehingga
setiap atribut yang dimasukkan pada diagram kelas memiliki alasan fungsional
yang jelas. Pada operasi inspeksi utama, entitas KONTAINER diposisikan sebagai
pusat penyimpanan identitas kontainer, hasil OCR, ringkasan manifes, hasil
pemindaian, keputusan sistem, dan lampiran bukti inspeksi. Entitas MANIFEST
menyimpan data manifes yang lebih lengkap, seperti asal dan tujuan, pihak
pengirim dan penerima, deskripsi muatan, berat, volume, dan dokumen terkait.
Sementara itu, entitas USER berfungsi sebagai konteks autentikasi dan
kepemilikan atau peninjauan data inspeksi.

\begin{figure}[htbp]
  \centering
  {\tikzumlset{font=\scriptsize}
  \resizebox{0.96\textwidth}{!}{%
  \begin{tikzpicture}
    \umlclass[x=0,y=0,width=4.1cm,fill=black!6]{Kontainer}{
      + nomorKontainer : string\\
      + statusInspeksi : enum\\
      + waktuTerakhir : datetime
    }{}

    \umlclass[x=-5.2,y=3.0,width=3.6cm]{User}{
      + idUser : string\\
      + peran : string
    }{}
    \umlclass[x=-5.2,y=0,width=3.9cm]{AuditEvent}{
      + aksi : string\\
      + waktu : datetime
    }{}
    \umlclass[x=5.2,y=0,width=3.9cm]{Manifest}{
      + nomorDokumen : string\\
      + ringkasanMuatan : string
    }{}
    \umlclass[x=0,y=-3.0,width=4.1cm]{Inspeksi}{
      + idInspeksi : string\\
      + waktuMulai : datetime\\
      + hasil : enum
    }{}
    \umlclass[x=-5.2,y=-6.0,width=3.9cm]{Deteksi}{
      + jenis : string\\
      + confidence : float
    }{}
    \umlclass[x=0,y=-6.0,width=4.1cm]{ArtefakVisual}{
      + uri : string\\
      + tipeMedia : string
    }{}
    \umlclass[x=5.2,y=-6.0,width=3.9cm]{Peristiwa}{
      + tipe : string\\
      + waktu : datetime
    }{}

    \umluniassociate[geometry=--, mult1={1}, mult2={0..*}]{User}{AuditEvent}
    \umluniassociate[geometry=--, mult1={0..*}, mult2={1}]{AuditEvent}{Kontainer}
    \umlassociate[geometry=--, mult1={1}, mult2={0..1}]{Kontainer}{Manifest}
    \umlcompose[geometry=--, mult1={1}, mult2={0..*}]{Kontainer}{Inspeksi}
    \umlcompose[geometry=--, mult1={1}, mult2={0..*}]{Inspeksi}{Deteksi}
    \umlcompose[geometry=--, mult1={1}, mult2={0..*}]{Inspeksi}{ArtefakVisual}
    \umluniassociate[geometry=--, mult1={1}, mult2={0..*}]{Inspeksi}{Peristiwa}
  \end{tikzpicture}}}
  \caption[Diagram kelas UML model data berpusat pada kontainer]{Diagram Kelas UML Model Data Berpusat pada Kontainer}
  \label{fig:er_diagram}
\end{figure}

\FloatBarrier

Gambar~\ref{fig:er_diagram} menunjukkan model data konseptual dalam bentuk
diagram kelas UML pada tingkat sasaran. Pusat model berada pada kelas
Kontainer karena identitas kontainer
menjadi jangkar untuk konteks manifes, riwayat inspeksi, bukti visual, peristiwa
operasional, dan audit. Entitas INSPEKSI pada model ini dibaca sebagai entitas
pendukung yang membentuk riwayat pemeriksaan kontainer, bukan sebagai pusat data
utama. Dengan demikian, gambar ini perlu dibaca sebagai representasi konseptual
pengembangan sistem, bukan sebagai cerminan satu-per-satu dari seluruh koleksi
yang aktif pada implementasi akhir.

\subsubsection{Mekanisme Komunikasi Data}

Komunikasi antarkomponen menggunakan dua bentuk utama, yaitu pertukaran data terstruktur melalui HTTP dan penyiaran video langsung melalui \textit{WebSocket}. Setiap komunikasi membawa \textit{metadata} identifikasi, waktu, dan data yang diperlukan dalam proses inspeksi.

Sistem menggunakan tiga kategori komunikasi utama untuk koordinasi
antarkomponen, yaitu pengiriman hasil OCR dari \textit{edge} ke layanan API
untuk identifikasi kontainer, pengiriman hasil pemindaian kerusakan, barang
berisiko, dan kategori muatan dari \textit{edge} ke layanan API, serta
penyiaran \textit{frame} beranotasi dari \textit{edge} ke aplikasi web untuk
video langsung. Seluruh komunikasi tersebut mengikuti prinsip struktur
mengikuti standar agar \textit{metadata} dan isi data tetap konsisten, kontrak layanan
yang jelas agar format dan semantik pertukaran data dipahami oleh pengirim dan
penerima, pemisahan fungsi antara jalur data transaksional dan jalur video
langsung, serta keterkaitan konteks agar hasil dari beberapa aliran inspeksi
dapat dihubungkan pada entitas kontainer yang sama.

\section{Alur Informasi dan Integrasi Hasil Inspeksi}

Bagian ini merangkum alur integrasi hasil identifikasi, hasil pemindaian, dan
penyajian informasi inspeksi kontainer. Fokus pembahasan diarahkan pada
hubungan antara komponen \textit{edge}, layanan API, penyimpanan data, aplikasi
web, dan kesiapan antarmuka terhadap sistem eksternal.

\subsection{Strategi Integrasi dengan Sistem Eksternal}

Sistem menggunakan pendekatan integrasi yang menempatkan layanan API sebagai
titik pertukaran data terstruktur, sementara lapisan integrasi menjaga
kesesuaian format dan semantik data dengan sistem eksternal. Pada realisasi
sistem yang dibangun, fokus integrasi berada pada konsistensi alur data
internal antara komponen \textit{edge}, layanan API, penyimpanan data, dan
aplikasi web. Untuk kebutuhan integrasi dengan ekosistem pelabuhan yang lebih
luas, desain ini menyiapkan antarmuka yang dapat diperluas tanpa mengubah
logika aplikasi inti.

Karakteristik pendekatan integrasi ini mencakup komunikasi data terstruktur
melalui antarmuka layanan yang jelas, pemisahan logika proses internal dari
detail komunikasi sistem eksternal, dukungan transformasi data agar model
internal dapat dipetakan ke kebutuhan pihak lain, serta titik perluasan bertahap
ketika integrasi dengan sistem eksternal direalisasikan.

\subsubsection{Pola Adapter untuk Sistem Eksternal}

Untuk mengisolasi logika aplikasi inti dari detail implementasi integrasi, rancangan sistem menggunakan \textbf{pola \textit{adapter}}. Setiap sistem eksternal direncanakan memiliki \textit{adapter} khusus yang mengimplementasikan antarmuka standar, sehingga penambahan atau penggantian sistem eksternal dapat dilakukan tanpa mengubah kode aplikasi inti.

Keuntungan pola \textit{adapter} terletak pada pemisahan tanggung jawab karena
logika aplikasi inti tidak bergantung pada detail implementasi sistem eksternal,
kemudahan substitusi saat sistem eksternal perlu diganti atau dinaikkan
versinya, kemudahan pengujian tanpa ketergantungan pada sistem eksternal
aktual, serta standardisasi antarmuka yang mempermudah pemahaman hubungan
integrasi.

Diagram komponen UML pada Gambar~\ref{fig:integration_architecture} disusun dengan
mengambil aktor dan sistem eksternal dari diagram konteks, kemudian
menempatkannya di luar batas sistem inti. Aliran yang masuk dan keluar dari
sistem dipusatkan melalui antarmuka integrasi agar perubahan pada sistem
eksternal tidak langsung memengaruhi layanan inspeksi inti. Pada tahap
realisasi tugas akhir ini, bagian tersebut berfungsi sebagai rancangan kesiapan
perluasan, sedangkan integrasi yang telah ditunjukkan secara konkret berada
pada hubungan antara \textit{edge}, layanan API, penyimpanan data, dan aplikasi
web.

\begin{figure}[htbp]
  \centering
  \resizebox{0.98\textwidth}{!}{%
  \begin{tikzpicture}[
      font=\scriptsize,
      component/.style={rectangle, draw=black, rounded corners=1.5pt, fill=black!2,
        align=center, minimum width=2.92cm, minimum height=0.78cm, inner sep=4pt},
      componentwide/.style={component, minimum width=3.28cm},
      externalcomp/.style={component, minimum width=3.22cm, fill=black!1},
      internal/.style={-{Latex[length=2.25mm,width=1.7mm]}, thick, shorten >=3pt, shorten <=3pt},
      external/.style={-{Latex[length=2.25mm,width=1.7mm]}, thick, dashed, shorten >=3pt, shorten <=3pt},
      trunk/.style={thick, dashed},
      legendbox/.style={rectangle, draw=black, rounded corners=2pt, fill=white, inner sep=4pt},
      component glyph/.pic={
        \draw (0,0) rectangle (0.16,-0.24);
        \draw (-0.07,-0.055) rectangle (0.045,-0.105);
        \draw (-0.07,-0.155) rectangle (0.045,-0.205);
      }
    ]
    \node[component] (edge) at (-5.75,1.30) {Cargovision Edge};
    \node[component] (web) at (-5.75,-0.52) {Aplikasi Web};
    \node[componentwide] (core) at (-2.05,0.36) {Core API};
    \node[componentwide] (store) at (-2.05,-1.58) {MongoDB dan R2};
    \node[componentwide] (interface) at (1.70,0.36) {Integration Interface};

    \node[externalcomp] (tosAdapter) at (6.80,2.02) {Adapter TOS/PCS};
    \node[externalcomp] (hublaAdapter) at (6.80,0.78) {Adapter Inaportnet};
    \node[externalcomp] (customsAdapter) at (6.80,-0.46) {Adapter Bea Cukai};
    \node[externalcomp] (partnerAdapter) at (6.80,-1.70) {Adapter Mitra Logistik};

    \foreach \n in {edge,web,core,store,interface,tosAdapter,hublaAdapter,customsAdapter,partnerAdapter} {
      \pic at ($(\n.north east)+(-0.20,-0.14)$) {component glyph};
    }

    \coordinate (coreInTop) at ($(core.west)+(0,0.22)$);
    \coordinate (coreInBottom) at ($(core.west)+(0,-0.22)$);
    \coordinate (leftFan) at (-4.08,0.36);
    \draw[internal] (edge.east) -- (edge.east -| leftFan) |- (coreInTop);
    \draw[internal] (web.east) -- (web.east -| leftFan) |- (coreInBottom);
    \draw[internal] (core.south) -- (store.north);
    \draw[external] (core.east) -- (interface.west);

    \coordinate (split) at (4.35,0.36);
    \coordinate (fanTop) at (4.35,2.02);
    \coordinate (fanHubla) at (4.35,0.78);
    \coordinate (fanCustoms) at (4.35,-0.46);
    \coordinate (fanBottom) at (4.35,-1.70);
    \draw[trunk] (interface.east) -- (split);
    \draw[trunk] (fanTop) -- (fanBottom);
    \foreach \p in {split,fanTop,fanHubla,fanCustoms,fanBottom} {
      \fill (\p) circle (1.05pt);
    }
    \draw[external] (fanTop) -- (tosAdapter.west);
    \draw[external] (fanHubla) -- (hublaAdapter.west);
    \draw[external] (fanCustoms) -- (customsAdapter.west);
    \draw[external] (fanBottom) -- (partnerAdapter.west);

    \node[legendbox, minimum width=4.72cm, minimum height=1.05cm] at (-4.65,-2.88) {};
    \node[anchor=west, font=\scriptsize\bfseries] at (-6.78,-2.55) {Keterangan};
    \draw[internal] (-6.78,-2.87) -- (-6.03,-2.87);
    \node[anchor=west] at (-5.84,-2.87) {realisasi internal};
    \draw[external] (-6.78,-3.18) -- (-6.03,-3.18);
    \node[anchor=west] at (-5.84,-3.18) {kesiapan integrasi eksternal};
  \end{tikzpicture}}
  \caption{Diagram Komponen UML Arsitektur Integrasi Sasaran dengan Sistem Eksternal Pelabuhan}
  \label{fig:integration_architecture}
\end{figure}

\FloatBarrier

Gambar \ref{fig:integration_architecture} menunjukkan diagram komponen UML untuk arsitektur integrasi
sasaran sistem dengan ekosistem pelabuhan yang lebih luas. Pada implementasi
yang direalisasikan dalam tugas akhir ini, integrasi aktual masih dibatasi pada
komponen internal proyek dan dependensi pendukung seperti autentikasi pengguna
serta penyimpanan artefak visual. Oleh karena itu, lapisan
\textit{Integration Interface} dan \textit{Integration Service} perlu dibaca
sebagai rancangan kesiapan perluasan, bukan sebagai klaim realisasi integrasi
eksternal penuh.

Lapisan integrasi bertanggung jawab memetakan model data internal sistem ke
format data yang digunakan oleh sistem eksternal. Dalam rancangan ini,
transformasi data mencakup penyesuaian nilai kategori sesuai format sistem
eksternal, konversi satuan pengukuran bila diperlukan, penyesuaian format waktu
dan identifikasi, penanganan informasi yang tidak tersedia, serta pemetaan
informasi yang hanya dimiliki sistem eksternal tertentu. Lapisan integrasi juga
dirancang untuk mendukung pemetaan dua arah ketika aliran data dari sistem
internal ke sistem eksternal maupun sebaliknya direalisasikan.

Selain pemetaan data, rancangan lapisan integrasi juga menyiapkan mekanisme
ketahanan ketika komunikasi dengan sistem eksternal mengalami kegagalan.
Strategi yang disiapkan mencakup percobaan ulang dengan interval yang
meningkat, sifat idempoten agar permintaan dapat dicoba ulang dengan aman,
antrean cadangan untuk menyimpan permintaan yang gagal, serta mekanisme
pembatas untuk menghentikan sementara pengiriman saat tingkat kegagalan tinggi.
Rancangan pemantauan integrasi juga menekankan pelacakan metrik seperti tingkat
keberhasilan dan waktu respons, peringatan otomatis untuk kondisi anomali, serta
log audit atas permintaan dan respons integrasi sebagai bentuk kesiapan
kepatuhan.

\subsection{Antarmuka Layanan}

Sistem menyediakan antarmuka layanan standar untuk komunikasi dengan
aplikasi klien. Pada implementasi aktual, antarmuka ini terutama digunakan
untuk autentikasi pengguna, pengelolaan manifes, pengelolaan kontainer, serta
penerimaan hasil pemindaian dari komponen \textit{edge}. Dukungan terhadap
sistem eksternal tetap diposisikan sebagai kesiapan perluasan arsitektur.

Secara konseptual, antarmuka layanan dikelompokkan ke dalam layanan
autentikasi dan manajemen pengguna, manajemen kontainer, manajemen manifes,
manajemen inspeksi, manajemen deteksi, manajemen media, dan layanan integrasi
untuk kesiapan koordinasi dengan sistem eksternal pada tahap lanjutan. Setiap
kategori layanan memiliki kontrak antarmuka yang terdefinisi dengan jelas,
termasuk format permintaan, respons, dan penanganan kesalahan yang konsisten.

\subsection{Alur Proses Utama}

Sistem mendukung alur proses inspeksi menyeluruh yang dimulai dari inisiasi
inspeksi, dilanjutkan dengan akuisisi data, pengunggahan ke sistem terpusat,
analisis visual, validasi manual, generasi laporan, dan penyajian hasil kepada
pihak terkait. Integrasi dengan sistem eksternal diposisikan sebagai arah
pengembangan lanjutan agar hasil inspeksi dapat masuk ke ekosistem informasi
pelabuhan yang lebih luas.

Selain alur inspeksi, sistem juga mendukung pemantauan dan pengelolaan status
layanan. Pemantauan ini mencakup ketersediaan aliran video, status komponen
\textit{edge}, dan ketersediaan layanan untuk membantu operator mengenali
gangguan operasional serta menentukan tindak lanjut yang diperlukan. Dalam
konteks kegagalan komunikasi, komponen sistem dirancang agar gangguan pada satu
jalur tidak langsung mengganggu seluruh fungsi secara serempak. Pemisahan
jalur antara data terstruktur dan video langsung membantu menjaga
keberlangsungan pemantauan maupun pencatatan data sesuai kondisi layanan yang
tersedia.

\subsection{Prinsip Keamanan}

Keamanan sistem otomasi inspeksi kontainer dirancang dengan prinsip \textbf{pertahanan berlapis}, yaitu pemisahan beberapa area kontrol keamanan untuk melindungi data, sistem, dan akses pengguna. Desain keamanan diarahkan agar selaras dengan persyaratan regulasi dan praktik terbaik industri untuk pengelolaan informasi.

Lapisan keamanan konseptual dirancang mencakup keamanan \textit{boundary}
untuk mengisolasi sistem dari ancaman eksternal, keamanan aplikasi melalui
validasi masukan dan praktik pengkodean aman, keamanan data untuk melindungi
informasi saat disimpan maupun dipertukarkan, serta kontrol akses berbasis
peran dengan beberapa tingkat otorisasi. Dari sisi tata kelola, rancangan ini
juga menekankan jejak audit untuk penelusuran dan investigasi, pemantauan
kepatuhan secara berkelanjutan terhadap regulasi dan standar keamanan,
manajemen insiden untuk respons dan pemulihan sistem, serta manajemen akses
atas identitas pengguna, peran, dan hak akses.


\section{Pembagian Fungsi Antarkomponen}

Bagian ini menjawab tujuan keempat tugas akhir, yaitu merancang pembagian
fungsi antarkomponen berdasarkan karakter beban pemrosesan lokal, persistensi
data, dan penyajian informasi dalam operasi inspeksi kontainer.
Arsitektur penempatan menggambarkan bagaimana komponen sistem didistribusikan
untuk mendukung responsivitas pengolahan di lapangan sekaligus menjaga
konsistensi data pada layanan terpusat.

\subsection{Strategi Penempatan Komponen}

Sistem mengadopsi strategi penempatan tersebar yang membagi komponen berdasarkan
karakteristik beban kerja dan kebutuhan responsivitas. Komponen yang membutuhkan
pemrosesan dekat dengan sumber data ditempatkan pada perangkat \textit{edge},
sedangkan komponen yang membutuhkan penyimpanan terpusat, pengelolaan data, dan
penyajian informasi ditempatkan pada layanan API, basis data, dan aplikasi web.

Komponen yang memerlukan konsolidasi informasi, akses ke data historis, dan
penyediaan data kepada pengguna ditempatkan pada layanan terpusat. Penempatan
ini membuat kapasitas dapat disesuaikan dengan volume operasi, data
terkonsolidasi dapat disimpan secara berkelanjutan, autentikasi dan akses
pengguna dapat dikelola seragam, serta antarmuka data untuk aplikasi web
tersedia dari satu titik layanan.

Komponen yang memerlukan responsivitas pada tahap pemantauan awal ditempatkan
di lokasi lokal. Komponen ini harus terhubung langsung dengan kamera
\textit{RTSP}, menjalankan inferensi awal dan pengolahan video di dekat sumber
data, mengirimkan hasil inspeksi ke layanan API, dan menyiarkan video
beranotasi langsung ke aplikasi web.

Koordinasi antara komponen lokal dan layanan terpusat dilakukan melalui
pengiriman hasil identifikasi dan pemindaian dari \textit{edge} ke layanan API,
penyediaan data inspeksi dari layanan API ke aplikasi web, penyiaran video
langsung dari \textit{edge} ke aplikasi web, serta pemantauan ketersediaan
aliran video dan status layanan.

Diagram deployment UML pada Gambar~\ref{fig:deployment_architecture} disusun berdasarkan
tiga kriteria. Pertama, fungsi yang membutuhkan kedekatan dengan sumber data
ditempatkan pada \textit{edge}. Kedua, fungsi yang membutuhkan konsolidasi,
riwayat, dan kontrol akses ditempatkan pada layanan terpusat. Ketiga, artefak
visual ditempatkan pada penyimpanan objek karena karakteristik ukurannya
berbeda dari data transaksional. Kriteria tersebut menjadi dasar pembagian
antara komponen lokal dan komponen terpusat.

\begin{figure}[htbp]
  \centering
  \resizebox{0.96\textwidth}{!}{%
  \begin{tikzpicture}[
      font=\scriptsize,
      deploy/.style={rectangle, draw=black, rounded corners=2pt, align=center, fill=black!2},
      device/.style={deploy, minimum width=3.65cm, minimum height=6.20cm},
      network/.style={deploy, minimum width=2.90cm, minimum height=6.20cm},
      cloud/.style={deploy, minimum width=7.10cm, minimum height=6.20cm},
      artifact/.style={rectangle, draw=black, align=center, fill=white, text width=2.52cm, minimum width=2.52cm, minimum height=0.76cm, inner xsep=5pt, inner ysep=4pt},
      database/.style={cylinder, draw=black, shape border rotate=90, aspect=0.25, align=center, fill=white, minimum width=2.35cm, minimum height=0.96cm},
      line/.style={-{Latex[length=2.05mm,width=1.55mm]}, thick, shorten >=1pt, shorten <=1pt},
      bendline/.style={line, rounded corners=3pt},
      videoline/.style={-{Latex[length=2.05mm,width=1.55mm]}, thick, dashed, shorten >=1pt, shorten <=1pt},
      label/.style={fill=white, inner sep=1.5pt, align=center},
      legendbox/.style={rectangle, draw=black, rounded corners=2pt, fill=white, inner sep=4pt}
    ]
    \node[device] (edgeDevice) at (-5.75,0) {};
    \node[anchor=north, align=center, font=\scriptsize\bfseries] at (edgeDevice.north) {$\ll$device$\gg$\\Lokasi Inspeksi};
    \node[network] (networkNode) at (-1.55,0) {};
    \node[anchor=north, align=center, font=\scriptsize\bfseries] at (networkNode.north) {$\ll$network$\gg$\\LAN / Internet};
    \node[cloud] (cloudNode) at (4.05,0) {};
    \node[anchor=north, align=center, font=\scriptsize\bfseries] at (cloudNode.north) {$\ll$node$\gg$\\Layanan Terpusat};

    \node[artifact] (camera) at (-5.75,1.78) {$\ll$device$\gg$\\Kamera RTSP};
    \node[artifact] (edge) at (-5.75,0.28) {$\ll$execution environment$\gg$\\Cargovision Edge};
    \node[artifact] (localBuffer) at (-5.75,-1.22) {$\ll$artifact$\gg$\\Buffer lokal};

    \node[artifact] (http) at (-1.55,0.28) {$\ll$protocol$\gg$\\HTTP API};
    \node[artifact] (ws) at (-1.55,-2.24) {$\ll$protocol$\gg$\\WebSocket};

    \node[artifact] (api) at (2.05,0.28) {$\ll$execution environment$\gg$\\Layanan API};
    \node[database] (mongo) at (5.78,0.28) {$\ll$database$\gg$\\MongoDB};
    \node[artifact] (r2) at (2.05,-1.35) {$\ll$artifact$\gg$\\R2};
    \node[artifact] (web) at (5.78,-1.35) {$\ll$artifact$\gg$\\Aplikasi Web};

    \draw[line] (camera.south) -- (edge.north);
    \draw[line] (edge.south) -- (localBuffer.north);

    \draw[line] (edge.east) -- (http.west);
    \coordinate (edgeVideoOut) at ($(edge.east)+(0,-0.08)$);
    \coordinate (videoBranch) at (-3.45,-2.24);
    \draw[videoline, rounded corners=3pt] (edgeVideoOut) -- ++(0.52,0) |- (videoBranch) -- (ws.west);

    \draw[line] (http.east) -- (api.west);
    \coordinate (webVideoIn) at ($(web.south west)+(0.55,0)$);
    \draw[videoline, rounded corners=3pt] (ws.east) -- ++(1.12,0) -| (webVideoIn);

    \draw[line] (api.east) -- (mongo.west);
    \draw[line] (api.south) -- (r2.north);

    \draw[line] (r2.east) -- (web.west);

  \end{tikzpicture}}
  \caption{Diagram Deployment UML dengan Distribusi Komponen Terpusat dan \textit{Edge}}
  \label{fig:deployment_architecture}
\end{figure}

\FloatBarrier

Gambar \ref{fig:deployment_architecture} menunjukkan diagram deployment UML antara perangkat \textit{edge}, jaringan, dan layanan terpusat. Pada sisi \textit{edge}, sistem menjalankan layanan untuk akuisisi aliran video, inferensi, dan siaran video langsung. Pada sisi terpusat, layanan API menangani persistensi dan penyediaan data, MongoDB menyimpan data terstruktur, media penyimpanan objek menyimpan artefak visual, dan aplikasi web menyediakan antarmuka bagi pengguna. Arsitektur ini menempatkan pemrosesan yang sensitif terhadap latensi di dekat sumber data, sementara penyimpanan dan penyajian informasi dipusatkan untuk menjaga konsistensi.

\subsection{Arsitektur \textit{Edge} dan Akuisisi Data}

Sebagai bagian dari penempatan lokal, sistem mengimplementasikan arsitektur \textit{edge} yang berfokus pada akuisisi aliran video dan pemrosesan responsif di lapangan.

Gambar yang disajikan pada bagian ini diturunkan dari kebutuhan untuk
memisahkan fungsi yang sensitif terhadap latensi dari fungsi yang menuntut
konsolidasi data terpusat. Karena itu, ilustrasi penempatan komponen tidak
sekadar menunjukkan lokasi layanan, tetapi juga menjelaskan tahap bagaimana
data bergerak dari kamera, diproses di \textit{edge}, dikirim ke layanan API,
dan akhirnya disajikan pada aplikasi web.

\begin{figure}[htbp]
  \centering
  \resizebox{0.92\textwidth}{!}{%
  \begin{tikzpicture}[
      font=\small,
      group/.style={rectangle, draw=black, rounded corners, minimum width=10.8cm, minimum height=0.95cm, align=center},
      line/.style={-{Latex[length=2mm]}, thick},
	      flowlabel/.style={fill=white, inner sep=1.5pt, font=\small}
	    ]
	    \useasboundingbox (-8.4,-4.95) rectangle (8.4,0.55);
	    \node[group] (cloud) at (0,0) {\textbf{Cloud / Layanan Terpusat}\\Layanan API, MongoDB, R2, dan Aplikasi Web};
    \node[group] (network) at (0,-1.45) {\textbf{Jaringan}\\HTTP untuk data terstruktur dan kanal video langsung};
    \node[group] (local) at (0,-2.9) {\textbf{Pemrosesan Lokal}\\OCR dan deteksi pada komponen lokal};
    \node[group] (sensor) at (0,-4.35) {\textbf{Sumber Data Lapangan}\\Kamera RTSP dan aliran video inspeksi};

    \draw[line] (sensor.north) -- (local.south);
    \draw[line] (local.north) -- (network.south);
    \draw[line] (network.north) -- (cloud.south);
    \draw[line, dashed] (network.east) -- ++(1.45,0) |- (cloud.east);
    \node[flowlabel, anchor=west] at (6.95,-0.72) {video langsung};
  \end{tikzpicture}}
  \caption{Arsitektur \textit{Edge} di Lokasi Inspeksi}
  \label{fig:iot_sensor_architecture}
\end{figure}

\FloatBarrier

Gambar \ref{fig:iot_sensor_architecture} menunjukkan arsitektur \textit{edge}
yang diimplementasikan di lokasi inspeksi. Pada realisasi saat ini, sumber data
visual berasal dari kamera \textit{RTSP}, baik melalui kamera IP berbasis
aplikasi maupun aliran video simulatif untuk pengujian. Kamera
\textit{RTSP} menjadi perangkat akuisisi utama untuk identifikasi kontainer,
pemindaian kerusakan, pemindaian barang berisiko, dan klasifikasi muatan.
Layanan \textit{edge} bertindak sebagai unit pemrosesan lokal yang mengelola
aliran video dari kamera, menjalankan inferensi OCR dan deteksi, membentuk
hasil inspeksi dalam format data terstruktur, mengirimkan hasil identifikasi
dan hasil pemindaian ke layanan API, serta menyiarkan video beranotasi langsung
ke aplikasi web. Dalam infrastruktur komunikasi, layanan API berperan sebagai
komponen pusat yang menerima data hasil inspeksi dari \textit{edge},
menyimpannya, dan menyediakannya kepada aplikasi web, sedangkan aplikasi web
menerima video langsung dari \textit{edge} dan membaca data terstruktur dari
layanan API.

Dalam alur operasinya, kamera mengirimkan aliran video ke layanan
\textit{edge} untuk pemrosesan awal. Hasil identifikasi dan hasil pemindaian
dikirim ke layanan API untuk penyimpanan permanen, sedangkan video beranotasi
dikirim langsung ke aplikasi web. Arsitektur ini mendukung responsivitas
pemantauan awal melalui pemrosesan lokal sekaligus menjaga konsistensi data pada
layanan terpusat.


\section{Pemetaan Komponen terhadap Atribut Kualitas ISO 25010}
\label{sec:pemetaan-iso25010}

Desain arsitektur sistem otomasi inspeksi kontainer secara eksplisit mempertimbangkan atribut kualitas yang didefinisikan dalam standar ISO/IEC 25010. Tabel~\ref{tbl:iso25010_mapping} memetakan komponen utama arsitektur terhadap atribut kualitas yang didukungnya.

\begin{table}[ht]
  \centering
  \small
  \caption[Pemetaan komponen terhadap atribut kualitas]{Pemetaan Komponen Arsitektur terhadap Atribut Kualitas ISO 25010}
  \label{tbl:iso25010_mapping}
  \setlength{\tabcolsep}{4pt}
  \renewcommand{\arraystretch}{1.08}
  \begin{tabularx}{\textwidth}{>{\raggedright\arraybackslash}p{0.23\textwidth} >{\raggedright\arraybackslash}p{0.22\textwidth} >{\raggedright\arraybackslash}X}
    \toprule
    \textbf{Komponen} & \textbf{Atribut ISO 25010} & \textbf{Kontribusi} \\
    \midrule
    Arsitektur Berlapis & \textit{Maintainability} - \textit{Modularity} & Pemisahan tanggung jawab memudahkan modifikasi komponen \\
    Pemisahan Jalur Komunikasi & \textit{Performance Efficiency} - \textit{Time Behaviour} & Pemisahan jalur data terstruktur dan video mendukung responsivitas dan konsistensi \\
    Penempatan Tersebar & \textit{Performance Efficiency} - \textit{Time Behaviour} & Pemrosesan lokal diarahkan untuk mengurangi waktu respons pada pemantauan awal \\
    Pola Adapter Integrasi & \textit{Compatibility} - \textit{Interoperability} & Abstraksi sistem eksternal memudahkan pertukaran informasi ketika integrasi diperluas \\
    Jejak Audit & \textit{Security} - \textit{Accountability} & Pencatatan aktivitas mendukung penelusuran tindakan pada ruang lingkup sistem \\
    Kontrol Akses Berbasis Peran & \textit{Security} - \textit{Authenticity} & Verifikasi identitas dan otorisasi melindungi akses ke fungsi sensitif \\
    Validasi Data Berlapis & \textit{Security} - \textit{Integrity} & Pemeriksaan data di berbagai lapisan memastikan kebenaran informasi \\
    Model Data Standar & \textit{Portability} - \textit{Adaptability} & Struktur data konsisten memudahkan adaptasi sistem \\
    Antarmuka Layanan Standar & \textit{Compatibility} - \textit{Co-existence} & Kontrak layanan memungkinkan koeksistensi dengan sistem lain \\
    Komponen \textit{Edge} & \textit{Reliability} - \textit{Availability} & Pengolahan dekat sumber data mendukung kesinambungan fungsi pemantauan awal \\
    \bottomrule
  \end{tabularx}
\end{table}

\FloatBarrier

Pemetaan pada tabel tersebut menunjukkan bahwa keputusan desain tidak berdiri
sendiri, tetapi dikaitkan dengan atribut kualitas yang ingin dicapai. Dari sisi
\textit{maintainability}, arsitektur berlapis dengan pemisahan tanggung jawab
yang jelas memastikan perubahan pada satu komponen tidak memengaruhi komponen
lain. Penggunaan antarmuka standar juga memungkinkan penggantian
implementasi komponen tanpa mengubah komponen yang bergantung padanya.

Dari sisi \textit{interoperability}, pola \textit{adapter} untuk integrasi
sistem eksternal dan penggunaan kontrak layanan standar memfasilitasi
pertukaran informasi dengan berbagai sistem yang heterogen di ekosistem
pelabuhan ketika tahap integrasi lanjutan direalisasikan. Dari sisi
riwayat kontainer, pencatatan aktivitas pada lapisan sistem mendukung agar
tindakan penting dapat ditelusuri untuk keperluan investigasi, penyelesaian
sengketa, dan kepatuhan regulasi. Adapun dari sisi keandalan, desain
penempatan tersebar dengan pemisahan fungsi lokal dan terpusat diarahkan untuk
menjaga responsivitas pemantauan awal sekaligus mempertahankan konsistensi
pencatatan data.
Kontrak layanan yang jelas antarkomponen membantu meningkatkan keandalan
pertukaran informasi dan memudahkan pengendalian ketika terjadi gangguan
operasional.

\section{Kesimpulan Desain}

Bab ini telah menyajikan desain arsitektur sistem otomasi inspeksi kontainer yang dirancang untuk mengatasi permasalahan yang telah diidentifikasi di Bab III. Desain ini berfokus pada integrasi komponen, pemrosesan responsif, dan konsistensi alur data sebagai dasar realisasi sistem yang diimplementasikan.

\subsection{Keputusan Arsitektural Kunci}

Berikut adalah keputusan arsitektural utama beserta rasionalisasinya:

\begin{enumerate}
  \item Arsitektur berlapis (\textit{layered architecture}) digunakan dengan
  pemisahan tanggung jawab yang jelas antara lapisan presentasi, aplikasi, data,
  dan akuisisi lokal. Arsitektur ini dipilih karena dapat memudahkan
  pemeliharaan dengan mengisolasi perubahan pada satu lapisan tanpa memengaruhi
  lapisan lain, mendukung evolusi sistem secara bertahap tanpa mengganggu logika
  inti, serta memfasilitasi pengujian terpisah per lapisan untuk meningkatkan
  kualitas dan keandalan sistem.

  \item Pemrosesan tersebar dengan kapabilitas lokal digunakan dengan
  mengombinasikan pemrosesan lokal di titik akuisisi data dan pemrosesan
  terpusat untuk analisis yang memerlukan konsolidasi data. Pendekatan ini dipilih berdasarkan
  karakteristik lingkungan operasi pelabuhan. Operasi pemantauan inspeksi
  membutuhkan respons yang memadai pada tahap awal, sehingga pemrosesan lokal memungkinkan deteksi awal
  anomali kritis tanpa menunggu respons dari sistem terpusat. Konektivitas
  jaringan pelabuhan juga dapat berubah-ubah, sehingga pemrosesan dekat sumber
  data membantu menjaga fungsi yang membutuhkan respons cepat. Selain itu,
  volume data inspeksi visual relatif besar, sehingga pemrosesan awal lokal
  membantu mengurangi beban pengiriman data.

  \item Pemisahan jalur komunikasi digunakan untuk membedakan jalur data
  terstruktur dan video langsung. Pendekatan ini dipilih karena data hasil
  inspeksi membutuhkan persistensi dan kontrol akses yang berbeda dari aliran
  video waktu nyata. Dengan pemisahan tersebut, layanan API tidak dibebani oleh
  arus video berukuran besar, sementara aplikasi web tetap dapat menampilkan
  kondisi lapangan secara langsung tanpa mengganggu penyimpanan data
  transaksional. Pemisahan jalur juga memperjelas tanggung jawab antarkomponen.

  \item Pola \textit{adapter} untuk integrasi digunakan agar setiap sistem
  eksternal dapat diakses melalui \textit{adapter} khusus yang mengimplementasikan
  antarmuka standar. Pola ini dipilih sebagai arah kesiapan integrasi mengingat
  karakteristik ekosistem pelabuhan. Sistem eksternal seperti TOS, PCS, dan INSW
  memiliki antarmuka yang heterogen, sehingga \textit{adapter} menyediakan
  abstraksi yang menyeragamkan akses ketika integrasi tersebut direalisasikan.
  Sistem eksternal juga dapat berubah atau diganti tanpa koordinasi langsung
  dengan sistem inspeksi, sehingga \textit{adapter} membantu mengisolasi
  perubahan tersebut dari logika aplikasi inti. Pola ini juga memungkinkan
  pengembangan dan pengujian rancangan integrasi tanpa ketergantungan langsung
  pada sistem eksternal aktual.
\end{enumerate}

Keputusan tersebut diturunkan dari kebutuhan arsitektural yang telah
diidentifikasi pada Bab III. Untuk memetakan hubungan antara masalah,
kebutuhan, keputusan desain, dan dasar evaluasi, hubungan tersebut diringkas
pada Tabel~\ref{tbl:traceability-rancangan-arsitektur}.

\begin{table}[ht]
  \centering
  \footnotesize
  \caption[Matriks pemetaan rancangan arsitektur]{Matriks Pemetaan Rancangan Arsitektur}
  \label{tbl:traceability-rancangan-arsitektur}
  \setlength{\tabcolsep}{4pt}
  \renewcommand{\arraystretch}{1.12}
  \begin{tabularx}{\textwidth}{>{\raggedright\arraybackslash}p{0.18\textwidth} >{\raggedright\arraybackslash}X >{\raggedright\arraybackslash}X >{\raggedright\arraybackslash}X}
    \toprule
    \textbf{Kebutuhan} & \textbf{Keputusan Arsitektur} & \textbf{Alasan Desain} & \textbf{Dasar Evaluasi} \\
    \midrule
    Fragmentasi hasil inspeksi antarkomponen &
    Menjadikan entitas kontainer sebagai pusat agregasi hasil OCR, pemindaian, manifes, dan bukti visual &
    Seluruh hasil inspeksi membutuhkan identitas bersama agar dapat ditelusuri tanpa entri manual berulang &
    Propagasi nomor kontainer dan penyimpanan hasil pemindaian pada entitas yang sama. \\
    Kebutuhan pemantauan waktu nyata dan pencatatan transaksional yang berbeda karakter &
    Memisahkan jalur video langsung dari jalur data terstruktur &
    Video menuntut latensi rendah, sedangkan data inspeksi membutuhkan validasi, persistensi, dan kontrol akses &
    Video beranotasi diterima melalui \textit{WebSocket}, sementara hasil inspeksi masuk melalui layanan API. \\
    Volume artefak visual dan data transaksional memiliki pola penyimpanan berbeda &
    Memisahkan basis data terstruktur dari penyimpanan objek untuk artefak visual &
    Data operasional membutuhkan query terstruktur, sedangkan citra hasil inspeksi lebih tepat disimpan sebagai objek &
    Dokumen kontainer menyimpan \textit{metadata} dan URL bukti, sedangkan citra tersimpan pada Cloudflare R2. \\
    Perubahan dan perluasan komponen harus dapat dilakukan bertahap &
    Menggunakan arsitektur berlapis dan kontrak layanan antarkomponen &
    Pemisahan tanggung jawab mengurangi ketergantungan langsung antarbagian sistem &
    Komponen \textit{edge}, API, web, dan penyimpanan dapat diuji berdasarkan peran masing-masing. \\
    Kesiapan integrasi dengan ekosistem pelabuhan yang heterogen &
    Menyiapkan pola \textit{adapter} dan antarmuka integrasi eksternal &
    Sistem eksternal memiliki format dan siklus perubahan yang berbeda sehingga perlu diisolasi dari logika inti &
    Pada ruang lingkup tugas akhir, pola ini divalidasi sebagai rancangan sasaran, bukan klaim integrasi eksternal penuh. \\
    \bottomrule
  \end{tabularx}
\end{table}

\FloatBarrier

Setiap keputusan arsitektural juga mengandung trade-off. Tabel
\ref{tbl:tradeoff-keputusan-arsitektur} merangkum alternatif yang tidak
dipilih agar alasan desain dapat ditelusuri secara eksplisit.

\begin{table}[ht]
  \centering
  \footnotesize
  \caption[Trade-off keputusan arsitektural]{Trade-off Keputusan Arsitektural}
  \label{tbl:tradeoff-keputusan-arsitektur}
  \setlength{\tabcolsep}{4pt}
  \renewcommand{\arraystretch}{1.12}
  \begin{tabularx}{\textwidth}{>{\raggedright\arraybackslash}p{0.20\textwidth} >{\raggedright\arraybackslash}X >{\raggedright\arraybackslash}X >{\raggedright\arraybackslash}X}
    \toprule
    \textbf{Keputusan} & \textbf{Alternatif} & \textbf{Alasan Tidak Dipilih} & \textbf{Konsekuensi Desain} \\
    \midrule
    Pemrosesan awal di \textit{edge} &
    Seluruh inferensi dilakukan di layanan terpusat &
    Membebani jaringan dan meningkatkan ketergantungan terhadap koneksi pusat untuk fungsi yang membutuhkan respons cepat &
    Komponen lokal perlu dikelola, tetapi responsivitas pemantauan meningkat. \\
    Jalur video langsung melalui \textit{WebSocket} &
    Video dikirim melalui layanan API yang sama dengan data transaksional &
    Menyatukan video dan data persisten membuat layanan API menanggung beban yang tidak sesuai dengan fungsi utamanya &
    Jalur komunikasi bertambah, tetapi tanggung jawab pemantauan dan persistensi menjadi lebih jelas. \\
    Artefak visual disimpan di penyimpanan objek &
    Citra hasil inspeksi disimpan langsung di basis data utama &
    Basis data transaksional menjadi berat dan kurang efisien untuk media berukuran besar &
    Diperlukan pengelolaan URL artefak, tetapi struktur data inti tetap ringan. \\
    Pola integrasi berbasis \textit{adapter} &
    Logika integrasi eksternal ditanam langsung pada layanan inti &
    Perubahan sistem eksternal berisiko memengaruhi logika inspeksi utama &
    Lapisan integrasi bertambah, tetapi sistem lebih siap terhadap perubahan eksternal. \\
    \bottomrule
  \end{tabularx}
\end{table}

\FloatBarrier

\subsection{Batas Operasional dan Ruang Lingkup Desain}

Desain arsitektur ini membutuhkan dukungan operasional berupa administrator
sistem, operator inspeksi, validator hasil, dan pengelola operasional. Keempat
peran tersebut diperlukan untuk menjalankan akuisisi data, validasi hasil,
pemantauan layanan, pengendalian akses, dan penanganan gangguan. Kebutuhan ini
selaras dengan temuan lapangan bahwa layanan inspeksi tidak hanya bergantung
pada perangkat teknis, tetapi juga pada prosedur, tingkat layanan, dan tata
kelola insiden.

Ruang lingkup desain difokuskan pada inspeksi visual nonintrusif terhadap
kondisi eksterior kontainer, dokumentasi kondisi dalam bentuk citra dan data
inspeksi, penyajian hasil kepada pengguna, serta kesiapan integrasi dengan
sistem informasi pelabuhan. Desain ini tidak mencakup inspeksi interior,
perbaikan fisik, perhitungan biaya kerusakan, validasi kepatuhan dokumen
pengiriman, penentuan tanggung jawab bisnis atau hukum, maupun integrasi dengan
sistem asuransi.

Dengan batas tersebut, rancangan tetap diarahkan pada kapabilitas inti
arsitektur: pemisahan tanggung jawab komponen, pengaitan hasil inspeksi pada
entitas kontainer, penyimpanan bukti visual, dan kesiapan integrasi bertahap.
